\subsection{Cosmological Descriptive Resolution}
  \label{subsec:cosmo_resolution}

  The parametrization $a_\star(z) \sim \kappa\, c_\chi^2 / \ell(z)$
  requires specifying the descriptive resolution in the homogeneous regime.

  In a spatially homogeneous late-time universe, no local preferred length scale exists.
  Operational closure therefore fixes $\ell(z)$ to the maximal domain
  over which an update of the projected state can coherently propagate
  within a single characteristic update time.

  The only intrinsic timescale in the homogeneous regime is $\tau(z) \sim H(z)^{-1}$,
  which follows directly from $H = \dot{\chi}/\chi$.
  The maximal operational influence radius within one update step is $c\,\tau(z) = c/H(z)$.

  Minimality of the closure then requires
  \begin{equation}
   \label{eq:ell_cosmo}
    \ell(z) = \frac{c}{H(z)}.
  \end{equation}

  Substituting Eq.~\eqref{eq:ell_cosmo} into Eq.~\eqref{eq:astar_param} gives
  \begin{equation}
    \label{eq:astar_general}
    a_\star(z) \sim \kappa\, \frac{c_\chi^2}{c}\, H(z),
  \end{equation}
  and in projectable regimes $c_\chi \to c$, so that
  \begin{equation}
    \label{eq:astar_projectable}
    a_\star(z) \sim \kappa\, c H(z).
  \end{equation}

  At the present epoch this yields $a_\star(t_0) \sim \kappa\, c H_0$,
  establishing a direct structural link between the galactic
  saturation scale and the cosmological expansion rate.
